\documentclass{article}
\usepackage{amsmath,hyperref}
\begin{document}
\section*{Fixed joint limits and the reserve needed for a later stop}

\paragraph{Decision and scope.}
A velocity command can satisfy the next position and command-change
bounds while leaving no lawful way to stop in later updates. This
companion to PR 855 holds the configured position, speed, and slew
bounds fixed, runs the same source limiter at every update, and
changes the sampled position through an explicit scalar plant
model. No ROS parameter service changes the bounds. The model
is sampled forward motion with exact following of each applied
velocity output, not a verified physical robot or ROS deployment.

\paragraph{Exact viability calculation.}
Normalize the sample period to one. Let $x$ be current position,
$W$ the fixed upper bound, $p\ge0$ the last applied command, and
$r>0$ the maximum commanded speed change per sample. The controller
may emit a nonnegative command $u$, subject to the fixed speed cap
and $|u-p|\le r$, and the model updates $x'=x+u$, $p'=u$.
For current $p$ the minimum remaining motion through a stop is
\[
  B_r(p)=\sum_{j\ge1}\max(0,p-jr).
\]
Any law-abiding sequence has $j$th subsequent command at least
$\max(0,p-jr)$ by induction and nonnegativity. The fastest monotone
brake attains the pointwise lower bound through zero without
exceeding the old speed cap. Hence $W-x\ge B_r(p)$ is exact,
including equality, provided the command and following assumptions
hold. For a chosen first command the exact test is
$u+B_r(u)\le W-x$, in addition to the present speed/slew constraints.
The minimal brake gives a feasible next action throughout this set.
For $q=\lfloor p/r\rfloor$, $B_r(p)=qp-rq(q+1)/2$; the formula
also includes an exact final zero term when $p/r$ is integral.
When $u$ is an actual speed for a nonunit fixed interval, multiply
the summed travel by that interval and take $r$ in speed units.

\paragraph{Native comparison.}
From rest under $W=7,r=1$, first request the outputs 1, 2, 3.
All three are emitted and the sampled positions are 1, 3, 6.
At the third update neither configured bound is violated, but the
remaining room is 1 and $B_1(3)=3$. At the fourth update even a
one-step intersection is empty. The normalized clamp still passes
the interior request 1.5 unchanged. The separate exact guard
substitutes 2 for 3 at the third update and can continue with
controlled motion while preserving sufficient rest distance. This
is a state transition in the explicit ideal plant, not evidence of
real motion dynamics. A cold initialization using a fresh measured
speed may change the software slew center; it cannot retroactively
alter the applied-command history or meet an impossible braking law.

\paragraph{What carries over.}
The \href{https://doi.org/10.1109/100.580977}{Dynamic Window
Approach of Fox, Burgard and Thrun} already excludes velocities from
which it cannot safely stop. A matched direct finite
braking-admissibility filter ties; the PR is a source-linked
counterexample to treating one-step clamp normalization as a stop
certificate, not a new planning advantage. The useful handoff
interface binds the last actual applied command, the sampled position
or conservative room bound, and a verified following law at the
same cut. Command feedback is insufficient if invocation does not
establish application, or the real plant cannot follow the demanded
deceleration.

\paragraph{Verification.}
The pinned source and explicit plant run 2,960 configurations/policy
pairs. Direct exact finite stopping search checks 2,624 states.
Of 1,480 unguarded traces, 870 leave the viable set and 864
show a first direct violation by the 12-sample horizon. All matched
guarded traces preserve modeled stoppability. Three Lean lemmas
check the pointwise braking bound and reserve steps; the real-valued
sharpness argument above is separate. This is a scoped R6 input,
not physical clearance or integration acceptance.
\end{document}
